\newpage
\section{The Propers: Themes and Movement}
\zlabel{triptych:concise:themes:start}
God's children receive their life from him and learn to desire what he promises. Their hope includes seeing God, belonging to his people and loving others with a charity still needing completion. Two interpretations explain these demands. Augustine's account of holy desire and Aquinas's account of happiness explain the person's preparation for vision. Bede's account of the multitude and Chrysostom's account of peacemaking explain the identity and conduct of the people Christ gathers. The prayers ask for the help both require.

\subsection{Rejoicing with the saints while asking for mercy}
The Entrance begins \textit{Gaudeamus omnes in Domino}, ``Let us all rejoice in the Lord.'' It celebrates the saints with angelic praise of God's Son. The congregation's joy concerns persons whose salvation comes from God. The Collect then asks for mercy with their intercession. Those celebrating the saints still need divine help; their petition belongs beside the thanksgiving with which they entered.

Revelation's crowd likewise attributes salvation to God and the Lamb. Its members come from nations, tribes, peoples and languages. Bede explains their praise as confession of gifts received from God in the Church's members.\footnote{Bede, \work{Explanation of the Apocalypse}, I, on 7:10--12, pp. 52--53.}

\subsection{Cleansed people who seek God's face}
The First Reading retains the number 144,000, omits the tribal counts and then describes a multitude beyond counting. Bede understands the number symbolically through apostolic fullness and the faith preached throughout the world. He permits both a resumed account of victorious glory and a succession from Israel to the Gentiles. The vision's breadth remains compatible with the suffering it describes. Bede expressly includes the whole Church in cleansing by the Lamb's blood, beyond those who shed their own blood in martyrdom.\footnote{Bede, \work{Explanation}, I, on 7:4,9,14, pp. 46,52--53.}

The psalm asks who can ascend God's mountain and stand in his holy place. Innocent hands and a pure heart describe action and desire together. Aquinas distinguishes approaching from abiding at the destination. He also permits blessing and mercy to explain how innocence becomes possible. Read beside the Lamb's cleansing, the psalm's demand for purity can thus be understood as a life made possible by received mercy.\footnote{Aquinas, \work{Super Psalmos} 23, nn. 2--4, paragraphs 86989--86991.}

First John states the problem of that life with particular clarity: God's children already belong to him, yet their future likeness remains to be disclosed. Hope leads them to purify themselves. Augustine explains the waiting through desire that enlarges the recipient's capacity for God's gift. He also explains willing purification through God's help: the person really consents and acts. Bede's cleansed multitude answers who stands before God; Augustine answers how the person hoping to stand there grows toward the promised sight.\footnote{Augustine, \work{Tractates on First John} IV.4--7.}

\newpage
\subsection{Learning from Christ and making peace}
Before the Gospel, Christ's invitation addresses the burdened in Matthew 11:28. In the following verses he invites them to learn his meekness and humility. Chrysostom explains how that learning makes the demanding yoke restorative. The humility he describes also appears at the beginning of his explanation of the Beatitudes: poverty of spirit concerns a chosen disposition, open to disciples of different social ranks.\footnote{Chrysostom, \work{Homilies on Matthew} XXXVIII.3; XV.2--3. Only Matthew 11:28 is appointed as the acclamation.}

Jesus teaches on a mountain before disciples, with the crowds in context. His blessings name mourning, hunger for justice and persecution alongside mercy, purity and peace. Augustine understands them as a course of conversion beginning with humility. He follows a text that places meekness before mourning; the U.S. proclaimed Gospel reverses those two. Aquinas distinguishes the acts preparing a person for happiness from completed happiness in the sight of God. The promises concern a fulfillment already begun and still awaited.\footnote{Augustine, \work{Our Lord's Sermon on the Mount} I.1--4, §§1--12; Aquinas, \work{Super Matthaeum}, cap. 5, lectio 2, paragraph 87340.}

Chrysostom explains peacemaking through the Son's reconciliation of those divided. The disciple must help other people become reconciled as well as refrain from personal quarrels. He then explains why Jesus immediately blesses persecution for justice: agreement can require abandoning virtue or help owed to others. The suffering blessed in the following direct address involves false accusations for Christ. A Christian seeking peace therefore still has to judge the conduct that produces agreement or opposition.\footnote{Chrysostom, \work{Matthew} XV.7--9. The Gospel ends at Matthew 5:12a; its following reference to prophets belongs to the wider context.}

\subsection{Intercession, Communion and charity's completion}
At the preparation of the gifts, the Prayer over the Offerings concerns the saints' continuing solicitude for those seeking salvation. Their completed life remains a life of concern for others. That intercession can be understood beside Chrysostom's account of active reconciliation: charity benefits persons who still need help. It also qualifies an account of purified desire, since desire fulfilled in God remains compatible with love for others.

The proper Preface names Jerusalem as mother and speaks of pilgrimage, the saints' help and their example. At Communion the antiphon repeats three blessings from Matthew 5:8--10: the pure heart, peacemaking and persecution for justice, with the promises of vision, sonship and kingdom. Augustine's account explains desire ordered toward God; Chrysostom's explains reconciliation among people. Hearing these promises together allows purity and charity toward others to be considered together, with endurance when fidelity meets hostility.

The Prayer after Communion asks concerning charity's perfection and arrival at the heavenly banquet. Present reception is followed by petition for further grace. The congregation hopes for the saints' common joy while still needing growth and perseverance. If the solemn blessing is selected, its petitions include intercession, example, service and final joy. Its concern for service belongs with the love already described in the readings and prayers; its actual selection remains optional.
\zlabel{triptych:concise:themes:end}
